% Part II: the Patient Catalog, the patient's page, the journey, the report, the detail, the
% encounters. Sources: docs/user-guide.md ("The Patient Catalog and the patient page",
% "Encounters and recordings", "Searching"); medorama/src/categories.rs; medorama/web/index.html,
% js/views/patients.js, patient.js, encounters.js, js/patient/*.js.
\chapter{Patients}
\label{chap:patients}

This chapter describes how to find a patient and read everything known about them. It is
written for clinicians, researchers and data managers. The \ui{Patient Catalog} tab lists every
patient of the catalog. A click on a patient opens the patient's page, which shows everything
known about the patient on one time line: who they are, their encounters, the Patient Journey, a
report of every entry and the detail of the entry chosen. The \ui{Encounters} table, beside the
list of patients, lists the encounters with their bed stays and recordings. How the patients'
data are brought in is the subject of \cref{part:data}.

\section{Concepts}
\label{sec:patients-concepts}

\begin{table}[htbp]
\caption{What the patient pages show.}
\label{tab:patients-concepts}
\begin{tabularx}{\linewidth}{@{}P{2.4cm}L@{}}
\toprule
{\tabhead Term} & {\tabhead Meaning} \\
\midrule
Patient & A person, known by a patient ID (\code{P900001}). The patient's identifiers (name,
  MRN, date of birth) are kept apart from the clinical data, in the catalog's identifier
  tables. \\
Encounter & One hospital stay or visit of a patient, known by an encounter ID
  (\code{E9100002}). \\
Bed stay & Where the patient was during an encounter: the unit, the bed, and the start and end
  times. \\
Item & One piece of a patient's data (a laboratory result, a vital sign, a drug, a code, a
  note, a report, an ECG, an image), belonging to a patient or an encounter and to a
  \emph{category} (\cref{tab:patients-categories}). \\
Recording & A waveform recording, such as ICU monitoring, a Holter ECG, an EEG or a sleep
  study. It is a WFDB record (single- or multi-segment), a folder of WFDB records read as one
  session, a WFDB collection, an EDF, EDF+, BDF or BDF+ file, or a file that an installed
  add-on reads. Each recording is linked to an encounter of its patient. \\
Entry & One line of a patient's report: an item, a panel of results taken together, a day of
  vital signs, or a recording. \\
\bottomrule
\end{tabularx}
\end{table}

\section{The Patient Catalog}
\label{sec:patients-catalog}

\screenshot{patients.png}{The Patient Catalog: every patient, latest encounter first, with the
kinds of data each has.}{fig:patients-catalog}

The list shows the patients, the one with the latest encounter first, a hundred to a page
(\ui{$\langle$ Previous}, \ui{Next $\rangle$}). Its columns are \ui{Patient} (the
name\phimark{} for users who see identifiers, otherwise the patient ID), \ui{Patient ID},
\ui{MRN}\phimark, \ui{Sex}, \ui{Age}, \ui{Born}\phimark, \ui{Encounters} (their number),
\ui{Latest encounter}, \ui{Where} (its unit), \ui{Data} (the categories of data the patient
has, and the connected databases that hold data about the patient) and \ui{Items}. A click on a
row opens the patient's page. In Privacy Mode, patients are named by their pseudonyms, and the
identifying columns are not shown.

\paragraph{Searching.} The search line finds patients as you type. A word matches any of the
searchable fields, \code{field:value} looks in one field only, and quotes keep words with
spaces together. The fields are the following:
\begin{itemize}
\item \code{patient}, \code{encounter}, \code{sex}, \code{unit}, and \code{has}, which names a
  category (\code{has:icu});
\item with the PHI permission and Privacy Mode off, also \code{name}, \code{first},
  \code{last}, \code{mrn} and \code{dob}\phimark.
\end{itemize}
In Privacy Mode, or without the PHI permission, a word never matches an identifier, and an
identifier field is an error. Patients are then found by their pseudonym (\code{MO-\dots}),
encounter, unit or sex. In Privacy Mode every search, filter and sort uses only what you are
shown: pseudonyms, tokens and placeholders, and each patient's dates as they are shown, so a
search for a real identifier or a real date finds nothing, and lists are ordered by the dates
you see. Such searches are written to the audit log (\code{privacy\_search}). \key{Enter} opens the patient when the search finds exactly one, and
\key{/} puts the focus in the search line from anywhere on the screen.

\paragraph{Filtering by kind of data.} After \ui{Filter by}, one button per category of data
(\cref{tab:patients-categories}) filters the list. A first click keeps the patients who have
data of that category, a second click keeps those who have none, and a third click clears the
filter. The filters combine (``Patients with ECG, without ICU''), and \ui{Remove filters} clears
them all.

\paragraph{Adding to a cohort.} \ui{Add to cohort\dots} adds every patient that the current
search and filters find to a cohort that the user may change: choose the cohort and click
\ui{Add \dots{} patients}. The patients whom the cohort's definition does not select are listed
in it as \emph{added manually} (\cref{chap:cohorts}).

\paragraph{A cohort's members.} A cohort's \ui{Open in Patient Catalog} opens the list limited
to its members. The chip \ui{Cohort: \dots} at the top says so, and its $\times$ shows every
patient again.

\paragraph{Encounters.} The buttons \ui{Patients} and \ui{Encounters} at the right switch
between the list of patients and the table of encounters (\cref{sec:patients-encounters}).

\section{The patient's page}
\label{sec:patients-page}

\screenshot{patient.png}{A patient's page: the patient bar, the encounters and conditions, the
Patient Journey, the report and the detail of an ECG.}{fig:patients-page}

From left to right, the page shows who the patient is, what happened (as a picture), and what
it was (as text). It has the following parts:
\begin{itemize}
\item the \textbf{search line}, across the page, which shows the patient as a chip ($\times$
  closes the page and returns to the Patient Catalog) and searches for another patient, with
  suggestions as you type (\key{$\uparrow$} \key{$\downarrow$} to choose, \key{Enter} to open,
  \key{Esc} to close);
\item the \textbf{patient bar}, across the page (\cref{sec:patients-bar});
\item the \textbf{patient column}, on the left: encounters, medications, conditions
  (\cref{sec:patients-column});
\item the \textbf{Patient Journey}, in the middle (\cref{sec:patients-journey});
\item the \textbf{report} and the \textbf{detail}, on the right (\cref{sec:patients-report},
  \cref{sec:patients-detail}).
\end{itemize}
The page's address holds the entry chosen, the encounter shown, and the report's search and
mode, so that the browser's Back and Forward buttons, reloading and sharing the address keep
them.

\subsection{The patient bar}
\label{sec:patients-bar}

The patient bar shows the following:
\begin{itemize}
\item the name\phimark{} (in Privacy Mode, \code{[Patient Name]}), and \ui{DEMO} for a patient
  of the demo;
\item the sex and age, with the date of birth\phimark{} for users who may see it (in Privacy
  Mode, ages above 89 are shown as \code{90+});
\item the patient ID and the pseudonym (in Privacy Mode, the pseudonym alone);
\item the MRNs\phimark, with their issuing authority when there are several;
\item the span of the data, from first to last, and the number of encounters (the catalog's, and
  the stays and visits of connected databases);
\item the allergies, in amber with a warning sign. The bar says ``No allergy recorded'' when a
  database holds allergies but none for the patient, and ``Allergies: not recorded'' when no
  database holds allergies;
\item when a connected database did not answer, its name followed by ``not available''. The
  page then shows the other sources, without what that database holds.
\end{itemize}

\subsection{The patient column}
\label{sec:patients-column}

\begin{description}
\item[\ui{Encounters}] Each encounter, newest first, with whether it is inpatient or
  outpatient, the units it went through (\code{Cardiology $\rightarrow$ ICU7}), its encounter
  ID, its length in days and its dates. A click shows only that encounter in the journey and
  the report, unfolded. A second click, or the encounter's chip above the journey, shows all of
  them again. The stays and visits held by connected databases follow, each with its source.
\item[\ui{Medications}] The drugs of the latest stay, with their dose and since or until when.
\item[\ui{Conditions}] The diagnoses, with their ICD-10 code and date.
\end{description}

\section{The Patient Journey}
\label{sec:patients-journey}

The journey is a picture of the patient's whole history. Time runs upward, with the newest
items on top, as in the report. There is one column per category, each with its color and its
name written vertically at the top.

\subsection{The categories}
\label{sec:patients-categories}

The categories are the clinical groupings a physician reads a chart by, after HL7 v2's
diagnostic service sections and FHIR's resources and observation categories, not the
databases the data came from. One database may fill several columns (an imaging index holds
imaging and echo), and several databases may fill one. Only the categories that some database
provides are shown. \Cref{tab:patients-categories} lists them in the journey's order, after the
\ui{Encounters} column, which draws the encounters as light bands behind everything.

\begin{xltabular}{\linewidth}{@{}P{2.4cm}P{1.4cm}L@{}}
\caption{The categories: the journey's columns, in order, with their keys and the words that
name them.}
\label{tab:patients-categories}\\
\toprule
{\tabhead Category} & {\tabhead Key} & {\tabhead Words that name it} \\
\midrule
\endfirsthead
\toprule
{\tabhead Category} & {\tabhead Key} & {\tabhead Words that name it} \\
\midrule
\endhead
\bottomrule
\endfoot
ICU & \code{icu} & icu, intensive care, critical care, monitoring, bedside monitoring, waveforms,
  bedmaster \\
Vitals & \code{vit} & vit, vital, vitals, vital signs, vital sign, flowsheet \\
Labs & \code{lab} & lab, labs, laboratory, chemistry, hematology, coagulation, blood gas, blood
  gases \\
Microbiology & \code{micro} & micro, microbiology, culture, cultures, sensitivity,
  sensitivities \\
Pathology & \code{path} & path, pathology, biopsy, biopsies, cytology, histology \\
ECG & \code{ecg} & ecg, ekg, 12-lead, 12 lead, electrocardiogram, muse \\
Holter & \code{holter} & holter, ambulatory ecg, event monitor, patch monitor \\
EEG & \code{eeg} & eeg, electroencephalography, electroencephalogram, neurophysiology, ieeg,
  seeg \\
Sleep & \code{sleep} & sleep, sleep study, sleep studies, polysomnography, psg \\
Echo & \code{echo} & echo, echocardiography, echocardiogram, tte, tee \\
Cath lab & \code{cath} & cath, cath lab, catheterization, cardiac catheterization, angiography,
  angiogram, coronary angiography, xa, ivus, oct, pci \\
Imaging & \code{img} & imaging, img, radiology, ct, mri, x-ray, xray, ultrasound, dicom, pacs \\
Procedures & \code{proc} & proc, procedure, procedures, surgery, operation, operations, cpt \\
Medications & \code{med} & med, meds, medication, medications, drug, drugs, mar, pharmacy \\
Transfusions & \code{blood} & blood, transfusion, transfusions, blood product, blood products,
  blood bank \\
Life support & \code{support} & support, life support, ventilation, ventilator, respiratory
  support, dialysis, crrt, ecmo \\
Conditions & \code{dx} & dx, condition, conditions, diagnosis, diagnoses, problem, problems,
  problem list, icd \\
Assessments & \code{obs} & obs, observation, observations, assessment, assessments, score,
  scores \\
Notes & \code{note} & note, notes, clinical note, clinical notes, document, documents, report,
  reports \\
Allergies & \code{allergy} & allergy, allergies, intolerance, intolerances (shown in the patient
  bar, not as a column) \\
\end{xltabular}

A category is named by its key, its label or any of its words, ignoring case. These names work
in a table's \code{category} column, in a table's file name, in a cohort's criteria
(\code{has holter}, \code{has sleep}) and in the report's search (\code{micro}). An item's
category is decided in this order:
\begin{enumerate}
\item the category that its row names, in a \code{category} column;
\item otherwise, the category chosen for its table when the table was imported;
\item otherwise, the category that its table's file name names (\file{sleep\_index.csv} is
  Sleep, \file{psg\_studies.csv} is Sleep, \file{eeg.csv} is EEG);
\item otherwise, a category from its domain and name: a drug is Medications, a condition
  Conditions, a procedure Procedures, a note Notes, an observation Assessments, a measurement
  named like a vital sign (such as \code{heart rate}, \code{spo2} or \code{map}) Vitals, and any
  other measurement Labs.
\end{enumerate}
\Cref{chap:importing} describes how tables are imported.

\subsection{What the columns show}
\label{sec:patients-columns}

\begin{itemize}
\item \textbf{Six-hour cells.} Every item fills its six-hour cell. Stays, drugs and recordings
  are drawn as bars.
\item \textbf{Thin columns.} Laboratory results, vital signs and medications are split into one
  thin column per test, sign or drug (the most frequent ones), so that a stay reads as a
  heatmap. Single categories get a narrow column, wide enough for their name written
  vertically.
\item \textbf{ICU is the stay, not the monitor.} The ICU column shows the whole time spent in an
  intensive-care unit, taken from the bed stays, as a light bar. Each recording made during the
  stay is a solid bar over it, from its first sample to its last. That span is known once the
  recording is indexed; before then, an hour is shown, dashed. Each bar opens its entry, however
  short. A unit counts as intensive care when its name says so (\code{ICU}, \code{ICU7},
  \code{MICU}, \code{CCU}, \code{ITU}, ``intensive'', ``critical care'' and so on).
\item \textbf{Holter, EEG and sleep recordings} each have their own column, with a bar as long
  as the recording.
\item \textbf{Abnormal values} are amber and \textbf{critical} ones orange, in every column. Red
  is reserved for alarms, and no category takes either color. A value is critical when its row
  flags it so (\code{HH}, \code{LL}, \code{C}, \code{critical}, \code{panic} and so on), and
  abnormal when it is flagged (\code{H}, \code{L}, \code{A}, \code{abnormal}, \code{positive},
  \code{resistant}, \code{detected} and so on) or lies outside its reference range.
\item \textbf{Folding.} Quiet stretches of more than twenty days shrink to a hatched band, so
  that busy days get the room.
\item \textbf{A long history scrolls} under fixed column names when years of care no longer
  fit. The journey opens at the top (the newest items), and choosing an entry in the report
  scrolls it into sight.
\end{itemize}

\subsection{Using the journey}
\label{sec:patients-journey-use}

\begin{itemize}
\item \textbf{Hover} over a mark to see what it is: the category, the test, how many items,
  whether they are abnormal or critical, and when.
\item \textbf{Click} a mark to open its entry in the report and the detail. For a cell, this is
  the newest entry of that column in that time. A click on an encounter's band shows only that
  encounter.
\item \textbf{Click a column's name} to hide that category from the journey and the report. Its
  name is then struck through, and a second click shows it again. The choice is remembered in
  this browser.
\item \textbf{Choose an encounter}, in the patient column or on its band, to show only that
  encounter, unfolded. Its chip above the journey ($\times$) shows all of them again (``All
  encounters'').
\item \ui{$-$} and \ui{+} zoom out and in.
\end{itemize}

\section{The report}
\label{sec:patients-report}

The report lists the patient's entries, newest first, in one of three modes:
\begin{description}
\item[\ui{Timeline}] The entries by day. Each day's heading gives the date and where the
  patient was: the encounter, the hospital day, and the unit and bed (\code{E9100002 $\cdot$
  hospital day 3 $\cdot$ ICU7 bed 12}). The laboratory results taken at one moment form one
  panel, and the vital signs form one entry per day.
\item[\ui{By category}] The entries grouped by category, in the journey's order.
\item[\ui{Trends}] Every measurement as a small chart (a sparkline), with its number of values,
  its last value and whether it is abnormal or critical. A click opens the entry of that value.
  A series is kept per name and unit: values in mg/dL and in mmol/L are two lines, not one.
\end{description}
Each entry shows its time, its category's color, its title, the number of abnormal values, a
summary, and an arrow when it opens a viewer or a file. A click (or \key{Enter}) chooses the
entry: the detail shows it, and the journey marks it.

\subsection{Searching a patient}
\label{sec:patients-search}

The report's search box (``Search everything about this patient: notes, results, codes\dots'')
looks through every text the patient has: the names of tests, drugs and codes, values, results,
notes and reports. The search shortens the report but does not change the journey. It reads the
patient's own items only, and in Privacy Mode their text as you are shown it: an entry is not
found by the patient's name, an identifier or a staff name that the page replaces. The words
found are highlighted, with a snippet of the text around them, and the number of entries found
in each category is shown above.

\screenshot{patient-search.png}{Searching a patient: the entries that mention
``creatinine'', the words highlighted.}{fig:patients-search}

\begin{itemize}
\item Words are alternatives: \code{ecg icu} finds the ECGs and the ICU monitoring.
\item A category's name finds all of that category: \code{labs}, \code{notes}, \code{ecg},
  \code{holter}, \code{sleep}, \code{micro}, \code{imaging} and so on
  (\cref{tab:patients-categories}).
\item Any part of a word matches (\code{lactat}), ignoring case. Quotes keep words together:
  \code{"two words"}.
\item \code{lactate > 4} and \code{hb<8} find values above or below a number (with \code{>},
  \code{<}, \code{>=}, \code{<=} or \code{=}).
\item A word found nowhere is matched to the closest word in the patient's text
  (\code{sepsys} finds \code{sepsis}), and the report says so (``\code{sepsys}: showing
  \code{sepsis}'').
\item The search does not understand negation: \code{no signs of sepsis} is found by
  \code{sepsis}.
\item The $\times$ in the box clears the search.
\end{itemize}

\begin{phibox}
The text of notes is searched and shown to users with the PHI permission who switched Privacy
Mode off. In Privacy Mode it is also searched and shown by default, with two protections: the
patient's own names and date of birth are replaced (\code{[Patient Name]}, \code{[Date of
Birth]}), and the text is marked with an amber warning triangle (``Free text: it may name the
patient (the de-identified view cannot remove names from it).''), because it may still name the
patient or others. An administrator can withhold notes in Privacy Mode instead, with \ui{The text
of notes and reports} in \ui{Admin} $\rightarrow$ \ui{Privacy Mode} (the key \code{notes};
\cref{sec:protecting-notes}). Notes are then found by their title only, and the detail says
``The text of notes is shown to users with the PHI permission''.
\end{phibox}

\section{The detail}
\label{sec:patients-detail}

The detail shows the entry chosen. Until one is chosen, it says ``Choose an entry, or a mark in
the journey''. Its heading gives the entry's title, when it was, where the patient was, and the
databases it came from. \ui{$\langle$} and \ui{$\rangle$} step to the previous and next entry of
the same kind. Below the heading, the detail shows, depending on the entry:
\begin{itemize}
\item \textbf{a panel of results}, as a table of \ui{Test}, \ui{Value}, \ui{Unit},
  \ui{Reference}, \ui{Previous} and \ui{This stay} (the test over the stay, as a small chart),
  with abnormal values flagged (\code{H}, \code{L}). A day of vital signs shows \ui{Sign},
  \ui{Last}, \ui{Unit}, \ui{Reference} and \ui{$\pm$12 h} (the twelve hours before and after);
\item \textbf{an item's fields}, as its row gave them, and its text (a note, a report);
\item \textbf{an attached file}:
  \begin{itemize}
  \item a 12-lead ECG (WFDB, GE MUSE XML, DICOM, EDF, JSON), drawn on paper by the ECG viewer
    itself in the detail, or a PNG, JPEG or PDF ECG shown as it is;
  \item a picture, a PDF or a video, shown in the detail;
  \item the buttons \ui{Open in ECG} or \ui{Open in PACS}, which open the file in that viewer's
    tab (\cref{chap:viewers-in-medorama}). ECGs open in the ECG viewer. Echo loops, angiography
    runs, images, DICOM series and folders, PDFs and videos of the Echo, Cath lab and Imaging
    categories open in the PACS viewer;
  \end{itemize}
\item \textbf{a recording} (\cref{sec:patients-recordings});
\item \textbf{the notes written around that time}, since notes belong to their encounter. The
  first one is open and the others open with a click; a second click opens the note as the
  entry.
\end{itemize}

\section{Recordings on the page}
\label{sec:patients-recordings}

\screenshot{patient-holter.png}{A Holter recording: its own column in the journey, and its heart
rate over the whole day in the detail.}{fig:patients-holter}

A recording is an entry of its patient's report, in its category: ICU, Holter, EEG or Sleep.
Its detail shows the recording at a glance:
\begin{itemize}
\item \textbf{an ICU recording} shows \ui{Recorded} (how much of its length holds samples),
  \ui{Waveforms} (its channels), \ui{Vital signs} (how many), \ui{Where} (unit, bed and
  format), and its main vital signs over its whole length, with its gaps;
\item \textbf{a recording with an ECG and no heart rate of its own}, such as a Holter ECG or a
  sleep study, shows \ui{Recorded}, \ui{Leads}, \ui{Heart rate} (the mean, in bpm) and
  \ui{Range}, and the heart rate over the whole recording.
\end{itemize}
A click on the chart (\ui{Whole recording}: ``click to open the viewer at that moment'') opens
VitalStream there, on a minute around the moment clicked. \ui{Open in VitalStream} opens the
recording at its start. Before the recording is indexed, the detail says ``Not indexed yet: its
overview appears once indexing is done.'' When no viewer is linked to the catalog, the detail
says so instead of offering the button.

\paragraph{Heart rate from the ECG.} A Holter ECG or a sleep study has no heart rate reported by
a monitor. For every recording that has an ECG lead and no heart rate of its own, the index job
derives one from the ECG; this is not a separate action. QRS complexes are found by a
Pan--Tompkins detector, an hour at a time. For each minute, the job then takes the median of the
beat-to-beat rates and their 5th to 95th percentile range. Minutes with fewer than 10 plausible
beats (30 to 250\,bpm) are left out (\cref{chap:actions}).

\paragraph{EEG recordings and sleep studies.} EEGs and sleep studies (polysomnography) usually
come as EDF, EDF+, BDF or BDF+ files, which Medorama reads without an add-on. They are
recordings like ICU monitoring and Holter ECGs, in the EEG and Sleep columns of the journey, and
they open in VitalStream. There, EDF+ annotations, such as sleep stages, apneas and arousals,
appear as events in the \ui{Annotations} lane, and an interrupted recording (EDF+D) shows a
break between its parts. The patient and recording fields of an EDF header can identify the
patient, so the viewers never show them; in Privacy Mode, EDF files and their annotations are
de-identified as described in \cref{sec:protecting-edf}. A folder of EDF files is not yet opened
as one session, and a file that holds only annotations (some scoring files do) is not yet shown.

\screenshot{patient-sleep.png}{The demo's overnight sleep study (an EDF+ file) on its patient's page: the Sleep column of the journey, and the study's entry in the report and the detail.}{fig:patients-sleep}

\Cref{chap:recordings} describes how recordings are brought in and linked to their encounters.

\section{Encounters and bed stays}
\label{sec:patients-encounters}

\ui{Encounters}, beside \ui{Patients} in the Patient Catalog, is the table of encounters. It has
one row per encounter, with a status dot for its recordings, the encounter, the patient, the
admission and discharge, its units, its number of recordings and its source. For users with the
PHI permission out of Privacy Mode, it also shows the MRN, the name and the date of
birth\phimark, and \ui{In the files} (\cref{sec:patients-in-files}).

\paragraph{Finding encounters.} The search line takes words or \code{field:value} terms
(\cref{tab:patients-fields}; the \ui{?} beside it lists the fields). \ui{Unit}, \ui{From},
\ui{To} and the status filter narrow the table. The search never depends on letter case, and a
value matches anywhere in the field: \code{name:smi} finds Smith, and \code{dob:1950-03} finds
March 1950. Every term must match, and an unknown field is an error whose message lists the
fields. In Privacy Mode every field holds what you are shown: \code{encounter} and
\code{patient} the pseudonyms, \code{file} the recordings' tokens (\code{mo-rec:12}),
\code{admit} and \code{discharge} the dates as shown. The same syntax serves the recordings of
\ui{Import} and the commands
\cmd{medorama list -{}-q}, \cmd{medorama recordings -{}-q} and \opt{-{}-search}.

\begin{table}[htbp]
\caption{The fields of the encounters' search (\phimark: only with the PHI permission, out of
Privacy Mode).}
\label{tab:patients-fields}
\begin{tabularx}{\linewidth}{@{}P{3.2cm}L@{}}
\toprule
{\tabhead Field} & {\tabhead Finds} \\
\midrule
\code{encounter}, \code{patient} & The encounter ID, the patient ID \\
\code{source} & The source that brought it (\code{source:"ICU pilot"}) \\
\code{unit}, \code{bed} & A unit or bed of its stays \\
\code{admit}, \code{discharge} & A date, in local time (\code{admit:2024-01}) \\
\code{disposition}, \code{sex}, \code{status} & The discharge disposition, the patient's sex,
  the recordings' status \\
\code{file} & The names of its linked recordings \\
\code{name}, \code{first}, \code{last}, \code{dob}, \code{mrn}, \code{csn}\phimark & The
  patient's identifiers, the encounter number \\
\code{file-name}, \code{file-id}, \code{file-dob}\phimark & What the linked recordings' files
  say about the patient \\
\bottomrule
\end{tabularx}
\end{table}

\paragraph{The status of recordings.} Each recording has a status, shown as a dot
(\cref{tab:patients-status}). An encounter shows the most urgent status among its recordings,
in this order: \code{missing}, \code{failed}, \code{indexing}, \code{queued}, \code{stale},
\code{new}, \code{ready}. An encounter without recordings has a dashed ring.

\begin{table}[htbp]
\caption{The status of a recording.}
\label{tab:patients-status}
\begin{tabularx}{\linewidth}{@{}P{1.6cm}P{2.6cm}P{3.8cm}L@{}}
\toprule
{\tabhead Status} & {\tabhead Filter} & {\tabhead Dot} & {\tabhead Means} \\
\midrule
\code{ready} & \ui{Indexed} & Green & Indexed, and the index matches the recording as it is
  now \\
\code{queued} & \ui{Queued} & Amber & Indexing is waiting for a worker \\
\code{indexing} & \ui{Indexing} & Amber, pulsing, with its percentage & Being indexed \\
\code{new} & \ui{Not indexed} & Gray & Never indexed \\
\code{stale} & \ui{Changed} & Purple ring & Changed since it was indexed \\
\code{failed} & \ui{Failed} & Amber & Indexing failed, or was cancelled; the error (or
  \emph{indexing was cancelled}) is shown. A cancelled recording is not indexed again on its
  own: \ui{Index now} indexes it \\
\code{missing} & \ui{Missing} & Amber & Not found at the last check \\
(none) & \ui{No recordings} & Dashed ring & An encounter without recordings \\
\bottomrule
\end{tabularx}
\end{table}

\paragraph{Selecting encounters.} The boxes at the start of the rows select encounters, and the
box in the header selects the whole page. \ui{Run actions on selected\dots} (for analysts and
administrators) runs actions on their recordings (\cref{chap:actions}). Administrators also have
\ui{Remove selected} and \ui{Remove matching\dots} (every encounter that the search finds). Both
first say what would go, and then remove the encounters with their items and links. The
recordings stay in the catalog, unlinked, and no file is touched.

\paragraph{The encounter panel.} A click on a row opens the encounter's panel at the side of the
page (\key{Esc} or $\times$ closes it) and adds the encounter to the user's History
(\cref{chap:page}). The panel shows:
\begin{itemize}
\item the patient (a link to the patient's page), the sex, the admission and discharge, the
  source and the disposition, and, for users with the PHI permission, the name, MRN, date of
  birth and encounter number\phimark;
\item \ui{Bed stays}, with \ui{Unit}, \ui{Bed}, \ui{From} and \ui{To};
\item \ui{Recordings}, each with its status, name, kind and format, how it was linked, its start
  and length, any error, and its buttons: \ui{View} (which opens it in VitalStream),
  \ui{Copy path}, \ui{Index now} (for analysts and administrators, when it is not indexed), and,
  for administrators, \ui{Remove}, and \ui{Unlink} or \ui{Forget manual link}
  (\cref{chap:recordings});
\item \ui{Items}, the encounter's items (all of them or by domain) and the patient's items
  outside encounters, with their time, domain, name, value, code and source;
\item for administrators, \ui{Remove this encounter from the catalog}.
\end{itemize}

\subsection{Identifiers in the recordings}
\label{sec:patients-in-files}

\begin{phibox}
Bedside monitors and EDF files often store patient identifiers inside the recordings
themselves. For users with the PHI permission out of Privacy Mode, the Encounters table shows
what the encounter's recordings say (\ui{In the files}), with a tick when it agrees with the
databases and a warning sign when it differs or names another patient. The encounter panel
shows these values beside the databases' values (``In the file: \dots''). The action
\code{identify} reads these identifiers into the catalog's database only (\cref{chap:actions}).
Recordings whose files disagree with the databases are listed for administrators in
\ui{Review} (\cref{chap:recordings}).
\end{phibox}

\section{Connected databases on the page}
\label{sec:patients-connected}

A patient's page shows the entries of every source in the same columns: the catalog's, and
those of the connected databases (an OMOP instance, a warehouse) whose persons are linked to the
patient. A connected database's stays and visits are listed after the catalog's encounters,
with the database's name. Its results, conditions, drugs and notes appear in the journey and the
report, and the detail names their source. These data are read from the database when the page
opens, and never copied. When the database does not answer, the patient bar says so.
\Cref{chap:connecting} describes how to connect databases.
